\section{Технический проект}

\subsection{Общая характеристика проектного решения}

Разрабатываемая программная система представляет собой web-сервис, предназначенный для размещения услуг, поиска исполнителей, выбора свободного времени и сопровождения взаимодействия между заказчиком и исполнителем после создания бронирования.

Система реализуется по клиент-серверной архитектуре. Клиентская часть представлена web-интерфейсом, который открывается в браузере пользователя и предоставляет доступ к функциям сервиса. Серверная часть выполняет обработку запросов, реализует бизнес-логику, проверяет права доступа, обеспечивает работу с данными и формирует HTML-ответы. В качестве постоянного хранилища используется реляционная база данных MySQL.

В основу проектного решения положено разделение системы на несколько связанных частей: клиентский интерфейс, серверное приложение, уровень доступа к данным, базу данных и файловое хранилище загружаемых материалов. Такое разделение позволяет сделать структуру проекта более понятной, выделить зоны ответственности отдельных компонентов и упростить дальнейшее сопровождение web-сервиса.

Важной особенностью решения является механизм календарной записи. Исполнитель задаёт интервалы работы по дням недели, количество недель действия расписания и перечень недоступных дат. На основе этих данных система формирует таблицу слотов, которые затем применяются при фильтрации услуг и оформлении заявок. Такой подход позволяет показывать пользователю только реально доступное время и снижает риск дублирования бронирований.

В процессе локальной разработки и тестирования система разворачивалась в среде Open Server Panel, включающей web-сервер Apache, интерпретатор PHP, СУБД MySQL и инструмент администрирования phpMyAdmin. Это позволило выполнять разработку, отладку и проверку структуры базы данных в едином локальном окружении.

На рисунке \ref{architecture:image} представлена общая схема организации решения задачи.

\begin{figure}[H]
	\center{\includegraphics[width=1\linewidth]{arch.png}}
	\caption{Общая схема организации решения задачи}
	\label{architecture:image}
\end{figure}

Как показано на схеме, пользователь обращается к сервису через браузер, который взаимодействует с серверной частью по протоколу HTTP. Серверная часть, реализованная на PHP, обрабатывает запросы, обращается к базе данных MySQL и возвращает клиенту сформированные страницы и результаты выполненных действий.

\subsection{Выбор средств и технологий реализации}

При выборе технологий для реализации web-сервиса учитывались простота развёртывания, совместимость с типовым web-окружением, поддержка серверной обработки форм, возможность работы с реляционной базой данных, удобство разработки интерфейса и достаточность выбранных средств для реализации функций каталога, расписания, бронирования, чата и отзывов.

В результате для реализации системы был выбран классический стек web-разработки, включающий HTML, CSS, JavaScript, PHP, MySQL и Apache. Данный стек позволяет реализовать необходимые элементы сервиса без использования тяжёлых внешних платформ и даёт возможность разворачивать систему как на локальном сервере, так и на стандартном хостинге.

\subsubsection{Состав применяемого технологического стека}

В процессе разработки web-сервиса используются следующие основные технологии:
\begin{itemize}
	\item HTML для формирования структуры web-страниц;
	\item CSS для стилизации и адаптивного оформления интерфейса;
	\item JavaScript для реализации клиентской интерактивности;
	\item PHP для обработки запросов и реализации серверной логики;
	\item MySQL для хранения данных;
	\item Apache для публикации и обработки HTTP-запросов;
	\item phpMyAdmin для администрирования и контроля структуры базы данных;
	\item Open Server Panel как локальная среда разработки и тестирования.
\end{itemize}

Каждая из перечисленных технологий используется для решения конкретного круга задач, возникающих при проектировании и реализации сервиса.

\subsubsection{Серверный язык программирования PHP}

В качестве основного языка программирования серверной части выбран PHP. Данный язык ориентирован на создание динамических web-приложений и широко применяется в системах, где требуется обработка HTML-форм, маршрутизация запросов, валидация данных и взаимодействие с базой данных.

В рамках проекта PHP используется для реализации следующих задач:
\begin{itemize}
	\item регистрация и авторизация пользователей;
	\item проверка ролей и прав доступа;
	\item создание, редактирование и отображение услуг;
	\item сохранение расписаний и недоступных дат;
	\item генерация свободных слотов записи;
	\item оформление бронирований и изменение их статусов;
	\item создание чатов и отправка сообщений;
	\item сохранение отзывов и избранных услуг;
	\item работа с загружаемыми изображениями.
\end{itemize}

Дополнительным преимуществом PHP является его хорошая совместимость с Apache и MySQL, что делает выбранный стек технологически согласованным и удобным для развёртывания.

\subsubsection{Клиентские web-технологии}

Клиентская часть сервиса реализована средствами HTML, CSS и JavaScript.

HTML используется для описания структуры страниц. С его помощью формируются карточки услуг, формы регистрации и авторизации, страницы профиля, элементы расписания, блоки переписки, списки бронирований и другие составные части интерфейса.

CSS применяется для задания единого визуального стиля системы. С использованием каскадных таблиц стилей реализуются оформление карточек услуг, адаптивные сетки, таблицы, кнопки, вкладки, элементы календаря, уведомления и модальные окна. Отдельное внимание при проектировании интерфейса уделялось адаптации страниц под мобильные устройства.

JavaScript используется для реализации интерактивного поведения пользовательского интерфейса. В рамках проекта он применяется для открытия и закрытия всплывающих элементов, клиентской работы с календарём, просмотра изображений, визуального обновления элементов интерфейса, обработки действий в чате и повышения удобства работы без полной перезагрузки страницы.

Выбор именно этих технологий обусловлен тем, что они являются стандартными средствами разработки web-интерфейсов и поддерживаются современными браузерами.

\subsubsection{Система управления базами данных MySQL}

В качестве системы управления базами данных выбрана MySQL. Данная СУБД обеспечивает хранение и выборку структурированных данных, поддержку внешних ключей, индексацию, сортировку, фильтрацию и выполнение SQL-запросов, необходимых для работы сервиса.

В рамках проекта база данных используется для хранения:
\begin{itemize}
	\item учётных записей пользователей;
	\item справочников городов и категорий;
	\item карточек услуг;
	\item изображений услуг;
	\item расписаний и недоступных дат;
	\item сгенерированных слотов;
	\item бронирований и их статусов;
	\item чатов и сообщений;
	\item изображений, прикреплённых к сообщениям;
	\item отзывов;
	\item избранных услуг пользователей.
\end{itemize}

Использование MySQL является обоснованным, поскольку данная СУБД широко применяется в web-разработке, хорошо совместима с PHP и предоставляет достаточные средства для реализации реляционной модели данных проектируемого сервиса.

\subsubsection{Средство администрирования базы данных phpMyAdmin}

В процессе разработки и тестирования базы данных использовался phpMyAdmin. Данное средство предоставляет web-интерфейс для просмотра структуры таблиц, выполнения SQL-запросов, редактирования данных, проверки связей между сущностями и экспорта содержимого базы данных.

Использование phpMyAdmin позволило:
\begin{itemize}
	\item визуально контролировать структуру базы данных;
	\item проверять корректность индексов и внешних ключей;
	\item анализировать содержимое таблиц на этапе тестирования;
	\item выполнять экспорт SQL-дампа базы данных;
	\item упрощать отладку работы серверной части при взаимодействии с MySQL.
\end{itemize}

Таким образом, phpMyAdmin использовался не как часть пользовательского сервиса, а как инструмент разработки, администрирования и проверки базы данных.

\subsubsection{Web-сервер Apache и локальная среда Open Server Panel}

В качестве web-сервера выбран Apache. Он обеспечивает приём HTTP-запросов, передачу их PHP-обработчикам и публикацию приложения в стандартной серверной среде.

На этапе локальной разработки использовалась среда Open Server Panel. Её применение позволило быстро развернуть локальное окружение, включающее Apache, PHP~8.3, MySQL~8.0 и phpMyAdmin, а также запускать и тестировать web-сервис без необходимости ручной настройки всех компонентов по отдельности.

Использование Open Server Panel было целесообразным по следующим причинам:
\begin{itemize}
	\item упрощённый запуск локального сервера;
	\item наличие готовых модулей Apache, PHP и MySQL;
	\item удобная работа с доменом проекта в локальной среде;
	\item встроенный доступ к phpMyAdmin;
	\item возможность быстрой отладки серверной логики и базы данных.
\end{itemize}

\subsection{Архитектурное построение web-сервиса}

Программная система построена по клиент-серверной модели. Пользователь взаимодействует с сервисом через браузер, который отправляет HTTP-запросы серверной части приложения. Серверная часть принимает и обрабатывает запросы, выполняет бизнес-логику, обращается к базе данных и файловому хранилищу, после чего формирует ответ для клиентской части.

С архитектурной точки зрения систему можно разделить на четыре основных уровня:
\begin{itemize}
	\item уровень пользовательского интерфейса;
	\item уровень серверной логики;
	\item уровень доступа к данным;
	\item уровень хранения данных и файлов.
\end{itemize}

На рисунке \ref{levels:image} представлена уровневая схема архитектуры web-сервиса.

\begin{figure}[H]
	\center{\includegraphics[width=1\linewidth]{levels.png}}
	\caption{Уровневая схема архитектуры web-сервиса}
	\label{levels:image}
\end{figure}

Уровень пользовательского интерфейса включает браузер пользователя, HTML-страницы, формы, таблицы стилей CSS и клиентские сценарии JavaScript. Через данный уровень пользователь выполняет основные действия в системе: просматривает каталог услуг, отправляет формы, выбирает доступные слоты, создаёт бронирования и взаимодействует с другими разделами сервиса.

Уровень серверной логики объединяет web-сервер Apache, PHP-приложение, маршрутизацию запросов, функциональные модули системы и общие компоненты проекта. На данном уровне выполняются обработка запросов, проверка пользовательских данных, работа с сессиями, выполнение бизнес-правил и формирование HTML-ответов.

Уровень доступа к данным обеспечивает взаимодействие серверной логики с базой данных. Для подключения к MySQL используется файл config/db.php, а выполнение SQL-запросов осуществляется через механизм PDO.

Уровень хранения данных и файлов включает базу данных MySQL и файловое хранилище проекта. В базе данных размещаются сведения о пользователях, услугах, бронированиях, слотах, сообщениях и отзывах. В файловом хранилище располагаются аватары пользователей, изображения услуг, изображения сообщений и статические ресурсы интерфейса.

Такое разделение позволяет отделить пользовательский интерфейс от серверной обработки, а серверную логику — от непосредственного хранения данных и файлов, что делает сопровождение и дальнейшее развитие web-сервиса более удобными.

\subsubsection{Компоненты уровня пользовательского интерфейса}

Уровень пользовательского интерфейса отвечает за отображение данных и взаимодействие пользователя с системой. Он включает HTML-страницы, формы ввода, элементы навигации, карточки услуг, кнопки действий, календарные элементы, модальные окна и клиентские сценарии JavaScript.

На рисунке \ref{clientcomp:image} представлена диаграмма компонентов уровня пользовательского интерфейса.

\begin{figure}[H]
	\center{\includegraphics[width=1\linewidth]{clientcomp.png}}
	\caption{Диаграмма компонентов уровня пользовательского интерфейса}
	\label{clientcomp:image}
\end{figure}

К основным компонентам данного уровня относятся:
\begin{itemize}
	\item браузер пользователя;
	\item HTML-страницы web-сервиса;
	\item формы регистрации, авторизации, редактирования профиля и создания услуг;
	\item карточки услуг и элементы каталога;
	\item элементы календаря и выбора доступных слотов;
	\item интерфейс чата;
	\item клиентские сценарии JavaScript;
	\item таблицы стилей CSS.
\end{itemize}

Данный уровень не выполняет непосредственного изменения данных в базе. Все действия пользователя передаются на серверную часть в виде GET- или POST-запросов, после чего сервер выполняет проверку данных и возвращает результат обработки.

\subsubsection{Компоненты уровня серверной логики}

Уровень серверной логики является центральной частью приложения. Он принимает запросы от web-сервера Apache, определяет необходимый обработчик, проверяет права доступа пользователя, выполняет валидацию данных и реализует основные бизнес-процессы системы.

На рисунке \ref{servercomp:image} представлена диаграмма компонентов уровня серверной логики.

\begin{figure}[H]
	\center{\includegraphics[width=1\linewidth]{servercomp.png}}
	\caption{Диаграмма компонентов уровня серверной логики}
	\label{servercomp:image}
\end{figure}

В состав серверной логики входят следующие основные компоненты:
\begin{itemize}
	\item главный файл index.php, обеспечивающий вход в приложение;
	\item модуль auth/, отвечающий за регистрацию, авторизацию и выход из системы;
	\item модуль services/, обеспечивающий просмотр каталога, карточек услуг и доступных слотов;
	\item модуль provider/, предназначенный для управления услугами, расписанием и заявками исполнителя;
	\item модуль booking/, отвечающий за создание бронирований и изменение статусов заявок;
	\item модуль chat/, обеспечивающий работу переписки по заявкам;
	\item модуль profile/, предназначенный для просмотра и редактирования профиля пользователя;
	\item модуль favorites/, реализующий добавление и удаление услуг из избранного;
	\item модуль reviews/, отвечающий за сохранение отзывов;
	\item файл includes/slot\_generator.php, выполняющий формирование доступных слотов;
	\item общие компоненты includes/, включающие шаблоны, сессии и пользовательские функции.
\end{itemize}

Функциональные модули серверного уровня используют общие компоненты системы. К ним относятся шаблоны верхней и нижней частей страниц, механизм работы с пользовательской сессией, функции проверки авторизации, перенаправления, валидации данных и вспомогательные операции, используемые в разных разделах приложения.

К основным алгоритмам серверного уровня относятся алгоритм формирования доступных слотов, алгоритм поиска услуг по выбранным параметрам, алгоритм создания бронирования и алгоритм добавления отзыва.

Алгоритм формирования слотов основывается на данных таблиц services, service\_schedules, service\_unavailable\_dates и slots. Исполнитель задаёт рабочие интервалы по дням недели и период действия расписания, после чего серверная часть с учётом длительности услуги формирует набор доступных временных интервалов. Если для услуги указаны исключённые даты, соответствующие интервалы не добавляются в список доступных слотов.

Алгоритм поиска услуг использует параметры фильтрации, выбранные пользователем: категорию, город и одну или несколько дат. Серверная часть анализирует карточки услуг и связанные с ними слоты, после чего возвращает только те предложения, по которым имеется доступное время в выбранные пользователем дни.

Алгоритм создания бронирования включает проверку авторизации заказчика, проверку существования услуги, контроль доступности выбранного слота, создание записи о бронировании и изменение состояния связанного временного интервала. Дополнительно для бронирования создаётся чат, в котором пользователи могут вести дальнейшую переписку по заявке.

Алгоритм добавления отзыва предусматривает проверку статуса бронирования и принадлежности текущего пользователя к участникам заявки. После успешной проверки серверная часть сохраняет отзыв и оценку в базе данных, а повторное добавление отзыва ограничивается на уровне бизнес-логики и структуры данных.

\subsubsection{Компоненты уровня доступа к данным}

Уровень доступа к данным обеспечивает взаимодействие серверной логики с базой данных MySQL. Для подключения к базе используется файл config/db.php, в котором задаются параметры соединения и создаётся объект подключения через PDO.

На рисунке \ref{datacomp:image} представлена диаграмма компонентов уровня доступа к данным.

\begin{figure}[H]
	\center{\includegraphics[width=1\linewidth]{datacomp.png}}
	\caption{Диаграмма компонентов уровня доступа к данным}
	\label{datacomp:image}
\end{figure}

К компонентам уровня доступа к данным относятся:
\begin{itemize}
	\item config/db.php, отвечающий за подключение к MySQL;
	\item механизм PDO, используемый для выполнения SQL-запросов;
	\item база данных MySQL;
	\item таблицы пользователей, услуг, расписаний, слотов, бронирований, сообщений, отзывов и избранных услуг;
	\item phpMyAdmin как средство администрирования и проверки структуры базы данных.
\end{itemize}

Серверная логика обращается к базе данных через общий файл config/db.php и механизм PDO. Это позволяет централизовать параметры подключения и исключить разрозненное обращение модулей приложения к базе данных.

\subsubsection{Компоненты уровня хранения файлов}

Отдельным уровнем в архитектуре выделяется хранение пользовательских файлов. Эти файлы физически размещаются в каталогах проекта, а в базе данных сохраняются только относительные пути к ним.

На рисунке \ref{filecomp:image} представлена диаграмма компонентов уровня хранения файлов.

\begin{figure}[H]
	\center{\includegraphics[width=1\linewidth]{filecomp.png}}
	\caption{Диаграмма компонентов уровня хранения файлов}
	\label{filecomp:image}
\end{figure}

К файловому уровню относятся:
\begin{itemize}
	\item каталог uploads/avatars для хранения изображений профилей пользователей;
	\item каталог uploads/services для хранения изображений услуг;
	\item каталог uploads/chat для хранения изображений, прикреплённых к сообщениям;
	\item каталог assets/ для хранения таблиц стилей, JavaScript-файлов и статических изображений интерфейса.
\end{itemize}

При загрузке пользовательских изображений серверная часть проверяет файл, сохраняет его в соответствующий каталог и записывает путь к нему в базу данных. При отображении страниц web-сервис получает путь из базы данных и передаёт его в HTML-разметку, после чего браузер загружает изображение как статический ресурс.

\subsubsection{Диаграмма последовательности обработки запроса}

На рисунке \ref{sequence:image} представлена диаграмма последовательности, иллюстрирующая взаимодействие компонентов системы при обработке запроса на создание бронирования.

\begin{figure}[H]
	\center{\includegraphics[width=1\linewidth]{sequence.png}}
	\caption{Диаграмма последовательности обработки запроса на создание бронирования}
	\label{sequence:image}
\end{figure}

Пользователь через браузер открывает страницу услуги и выбирает доступный временной слот. Браузер отправляет GET-запрос к services/view.php, после чего серверная часть получает из базы данных сведения об услуге и доступных слотах. После выбора слота выполняется переход к booking/create.php, где через подключение config/db.php / PDO запрашивается информация об услуге, исполнителе и выбранном слоте. Полученные данные используются для формирования страницы подтверждения бронирования.

После подтверждения записи браузер отправляет POST-запрос в booking/store.php. На этом этапе серверная часть проверяет роль пользователя, идентификатор услуги, идентификатор слота и комментарий заказчика. При возникновении ошибки пользователь возвращается на форму подтверждения с соответствующим сообщением. При успешной проверке начинается транзакция, в ходе которой выбранный слот блокируется запросом SELECT ... FOR UPDATE, создаётся запись о бронировании, создаётся чат, добавляются сообщения и обновляется статус слота на booked. После успешного выполнения всех операций транзакция фиксируется, а пользователь перенаправляется обратно на страницу услуги с сообщением об успешном создании заявки.

\subsubsection{Форматы данных}

Обмен данными между браузером пользователя и web-сервером Apache выполняется по протоколу HTTP. Для получения страниц и передачи параметров фильтрации используются GET-запросы, а для отправки форм регистрации, авторизации, бронирования, сообщений и отзывов — POST-запросы.

Данные обычных HTML-форм передаются в формате application/x-www-form-urlencoded. При загрузке изображений профиля, услуг и вложений сообщений используется формат multipart/form-data.

Результаты обработки запросов формируются серверной частью на PHP в виде HTML-страниц или перенаправлений. Данные, полученные из базы MySQL через PDO, используются PHP-скриптами при формировании интерфейса, отображаемого в браузере пользователя.

\subsection{Проект базы данных web-сервиса}

\subsubsection{ER-модель данных}

На основе анализа предметной области были выделены основные сущности системы и связи между ними. ER-модель базы данных отражает структуру хранения пользователей, услуг, расписаний, слотов, бронирований, сообщений, отзывов и избранных услуг.

К числу основных сущностей относятся:
\begin{itemize}
	\item пользователи;
	\item услуги;
	\item категории услуг;
	\item города;
	\item изображения услуг;
	\item расписания услуг;
	\item недоступные даты;
	\item слоты;
	\item бронирования;
	\item чаты;
	\item сообщения;
	\item изображения сообщений;
	\item отзывы;
	\item избранные услуги.
\end{itemize}

На рисунке \ref{er:image} представлена ER-диаграмма базы данных.

\begin{figure}[H]
	\center{\includegraphics[width=1\linewidth]{ermodel.png}}
	\caption{ER-модель базы данных}
	\label{er:image}
\end{figure}

Пользователь может выступать как заказчиком, так и исполнителем в зависимости от значения поля роли. Исполнитель связан с опубликованными услугами, а заказчик — с созданными бронированиями, отзывами и избранными услугами. Услуга связана с категорией, городом, изображениями, расписанием, недоступными датами и наборами слотов. Каждое бронирование связано с услугой, заказчиком, чатом и последующими сообщениями. После завершения взаимодействия по бронированию пользователи могут оставлять двусторонние отзывы.

\subsubsection{Нормализация данных}

Схема базы данных проектировалась с учётом требований нормализации. Основные таблицы содержат атомарные атрибуты, каждая сущность имеет первичный ключ, а повторяющиеся и справочные данные вынесены в отдельные таблицы. Связи между сущностями реализованы через внешние ключи.

Справочники категорий и городов отделены от таблиц пользователей и услуг, что позволяет избежать дублирования текстовых значений. Изображения услуг и сообщений вынесены в отдельные таблицы, поскольку одна услуга или одно сообщение могут быть связаны с несколькими файлами. Перечень недоступных дат хранится отдельно от расписаний, так как представляет собой независимое множество исключений. Слоты вынесены в отдельную таблицу, что позволяет явно хранить доступные интервалы и их текущее состояние.

Таким образом, схема данных соответствует требованиям третьей нормальной формы и обеспечивает согласованное хранение взаимосвязанных сущностей сервиса.

\subsubsection{Реляционная модель данных}

На основе ER-модели построена реляционная модель данных. В ней каждая сущность представлена отдельной таблицей, а связи между сущностями реализованы через первичные и внешние ключи.

На рисунке \ref{rel:image} представлена реляционная модель базы данных.

\begin{figure}[H]
	\center{\includegraphics[width=1\linewidth]{relmodel.png}}
	\caption{Реляционная модель базы данных}
	\label{rel:image}
\end{figure}

В реляционной модели основными являются таблицы users, services, bookings, slots, messages и reviews. Дополнительные таблицы обеспечивают хранение изображений, справочных значений, чатов, избранных услуг и параметров расписания.

\subsubsection{Описание структуры базы данных}

\setlength{\LTleft}{0pt}
\setlength{\LTright}{0pt}
\setlength{\tabcolsep}{3pt}
\renewcommand{\arraystretch}{1.1}

В таблице \ref{dbtables:table} приведён состав основных таблиц базы данных и их назначение.

\begin{longtable}{|>{\raggedright\arraybackslash}p{3.9cm}|>{\raggedright\arraybackslash}p{4.5cm}|>{\raggedright\arraybackslash}p{5.9cm}|}
	\caption{Основные таблицы базы данных web-сервиса}\label{dbtables:table} \\ \hline
	\centering Таблица & \centering Назначение & \centering\arraybackslash Основные связи \\ \hline
	\endfirsthead	
	\multicolumn{3}{l}{Продолжение таблицы \ref{dbtables:table}} \\ \hline
	\centering Таблица & \centering Назначение & \centering\arraybackslash Основные связи \\ \hline
	\endhead
	users & Хранение сведений о пользователях системы & Связана с услугами, бронированиями, сообщениями, отзывами и избранным \\ \hline
	categories & Справочник категорий услуг & Связана с таблицей services \\ \hline
	cities & Справочник городов & Связана с таблицами users и services \\ \hline
	services & Хранение карточек услуг исполнителей & Связана с users, categories, cities, bookings, service\_images, service\_schedules, service\_unavailable\_dates, slots \\ \hline
	service\_images & Хранение изображений услуг & Связана с services \\ \hline
	service\_schedules & Хранение регулярного расписания услуги & Связана с services \\ \hline
	service\_unavailable\newline\_dates & Хранение недоступных дат услуги & Связана с services \\ \hline
	slots & Хранение сформированных временных слотов & Связана с services и bookings \\ \hline
	bookings & Хранение заявок на бронирование & Связана с services, users, slots, chats, reviews \\ \hline
	chats & Хранение чатов по заявкам & Связана с bookings и messages \\ \hline
	messages & Хранение сообщений в чатах & Связана с chats, users и message\_images \\ \hline
	message\_images & Хранение изображений в сообщениях & Связана с messages \\ \hline
	reviews & Хранение отзывов и оценок пользователей & Связана с bookings и users \\ \hline
	favorites & Хранение избранных услуг пользователей & Связана с users и services \\ \hline
\end{longtable}

В состав таблицы users входят атрибуты, представленные в таблице \ref{users:table}.

\begin{longtable}{|>{\raggedright\arraybackslash}p{3.8cm}|>{\raggedright\arraybackslash}p{3.8cm}|>{\centering\arraybackslash}p{2.0cm}|>{\raggedright\arraybackslash}p{5.4cm}|}
	\caption{Атрибуты таблицы users}\label{users:table} \\ \hline
	\centering Поле & \centering Тип & \centering Обязат. & \centering\arraybackslash Описание \\ \hline
	\endfirsthead
	\multicolumn{4}{l}{Продолжение таблицы \ref{users:table}} \\ \hline
	\centering Поле & \centering Тип & \centering Обязат. & \centering\arraybackslash Описание \\ \hline
	\endhead
	id & INT & Да & Уникальный идентификатор пользователя \\ \hline
	email & VARCHAR(100) & Да & Адрес электронной почты для входа в систему \\ \hline
	password & VARCHAR(255) & Да & Хеш пароля пользователя \\ \hline
	full\_name & VARCHAR(100) & Да & Полное имя пользователя \\ \hline
	phone & VARCHAR(20) & Нет & Контактный номер телефона \\ \hline
	avatar & VARCHAR(255) & Нет & Путь к изображению профиля \\ \hline
	city\_id & INT & Нет & Ссылка на город пользователя \\ \hline
	about & TEXT & Нет & Дополнительная информация о пользователе \\ \hline
	role & ENUM(customer, provider) & Да & Роль пользователя в системе \\ \hline
	created\_at & TIMESTAMP & Да & Дата и время создания учётной записи \\ \hline
	updated\_at & TIMESTAMP & Да & Дата и время последнего обновления записи \\ \hline
\end{longtable}

В состав таблицы categories входят атрибуты, представленные в таблице \ref{categories:table}.

\begin{longtable}{|>{\raggedright\arraybackslash}p{3.8cm}|>{\raggedright\arraybackslash}p{3.8cm}|>{\centering\arraybackslash}p{2.0cm}|>{\raggedright\arraybackslash}p{5.4cm}|}
	\caption{Атрибуты таблицы categories}\label{categories:table} \\ \hline
	\centering Поле & \centering Тип & \centering Обязат. & \centering\arraybackslash Описание \\ \hline
	\endfirsthead
	\multicolumn{4}{l}{Продолжение таблицы \ref{categories:table}} \\ \hline
	\centering Поле & \centering Тип & \centering Обязат. & \centering\arraybackslash Описание \\ \hline
	\endhead
	id & INT & Да & Уникальный идентификатор категории \\ \hline
	name & VARCHAR(100) & Да & Наименование категории услуг \\ \hline
\end{longtable}

В состав таблицы cities входят атрибуты, представленные в таблице \ref{cities:table}.

\begin{longtable}{|>{\raggedright\arraybackslash}p{3.8cm}|>{\raggedright\arraybackslash}p{3.8cm}|>{\centering\arraybackslash}p{2.0cm}|>{\raggedright\arraybackslash}p{5.4cm}|}
	\caption{Атрибуты таблицы cities}\label{cities:table} \\ \hline
	\centering Поле & \centering Тип & \centering Обязат. & \centering\arraybackslash Описание \\ \hline
	\endfirsthead
	\multicolumn{4}{l}{Продолжение таблицы \ref{cities:table}} \\ \hline
	\centering Поле & \centering Тип & \centering Обязат. & \centering\arraybackslash Описание \\ \hline
	\endhead
	id & INT & Да & Уникальный идентификатор города \\ \hline
	name & VARCHAR(100) & Да & Наименование города \\ \hline
\end{longtable}

В состав таблицы services входят атрибуты, представленные в таблице \ref{services:table}.

\begin{longtable}{|>{\raggedright\arraybackslash}p{3.8cm}|>{\raggedright\arraybackslash}p{3.8cm}|>{\centering\arraybackslash}p{2.0cm}|>{\raggedright\arraybackslash}p{5.4cm}|}
	\caption{Атрибуты таблицы services}\label{services:table} \\ \hline
	\centering Поле & \centering Тип & \centering Обязат. & \centering\arraybackslash Описание \\ \hline
	\endfirsthead
	\multicolumn{4}{l}{Продолжение таблицы \ref{services:table}} \\ \hline
	\centering Поле & \centering Тип & \centering Обязат. & \centering\arraybackslash Описание \\ \hline
	\endhead
	id & INT & Да & Уникальный идентификатор услуги \\ \hline
	user\_id & INT & Да & Ссылка на исполнителя, разместившего услугу \\ \hline
	category\_id & INT & Да & Ссылка на категорию услуги \\ \hline
	city\_id & INT & Да & Ссылка на город оказания услуги \\ \hline
	title & VARCHAR(200) & Да & Название услуги \\ \hline
	description & TEXT & Нет & Подробное описание услуги \\ \hline
	price & DECIMAL(10,2) & Да & Стоимость услуги \\ \hline
	booking\_type & ENUM(timeslot, date\_only, flexible) & Да & Тип записи на услугу \\ \hline
	duration\_minutes & INT & Да & Длительность услуги в минутах \\ \hline
	status & ENUM(active, hidden, archived) & Да & Статус отображения услуги \\ \hline
	location\_text & VARCHAR(255) & Нет & Текстовое описание места оказания услуги \\ \hline
	created\_at & TIMESTAMP & Да & Дата и время создания услуги \\ \hline
	updated\_at & TIMESTAMP & Да & Дата и время последнего обновления услуги \\ \hline
\end{longtable}

В состав таблицы service\_images входят атрибуты, представленные в таблице \ref{serviceimages:table}.

\begin{longtable}{|>{\raggedright\arraybackslash}p{3.8cm}|>{\raggedright\arraybackslash}p{3.8cm}|>{\centering\arraybackslash}p{2.0cm}|>{\raggedright\arraybackslash}p{5.4cm}|}
	\caption{Атрибуты таблицы service\_images}\label{serviceimages:table} \\ \hline
	\centering Поле & \centering Тип & \centering Обязат. & \centering\arraybackslash Описание \\ \hline
	\endfirsthead
	\multicolumn{4}{l}{Продолжение таблицы \ref{serviceimages:table}} \\ \hline
	\centering Поле & \centering Тип & \centering Обязат. & \centering\arraybackslash Описание \\ \hline
	\endhead
	id & INT & Да & Уникальный идентификатор изображения услуги \\ \hline
	service\_id & INT & Да & Ссылка на услугу \\ \hline
	image\_path & VARCHAR(255) & Да & Путь к изображению услуги \\ \hline
	sort\_order & INT & Нет & Порядок отображения изображения \\ \hline
\end{longtable}

В состав таблицы service\_schedules входят атрибуты, представленные в таблице \ref{schedules:table}.

\begin{longtable}{|>{\raggedright\arraybackslash}p{3.8cm}|>{\raggedright\arraybackslash}p{3.8cm}|>{\centering\arraybackslash}p{2.0cm}|>{\raggedright\arraybackslash}p{5.4cm}|}
	\caption{Атрибуты таблицы service\_schedules}\label{schedules:table} \\ \hline
	\centering Поле & \centering Тип & \centering Обязат. & \centering\arraybackslash Описание \\ \hline
	\endfirsthead
	\multicolumn{4}{l}{Продолжение таблицы \ref{schedules:table}} \\ \hline
	\centering Поле & \centering Тип & \centering Обязат. & \centering\arraybackslash Описание \\ \hline
	\endhead
	id & INT & Да & Уникальный идентификатор записи расписания \\ \hline
	service\_id & INT & Да & Ссылка на услугу \\ \hline
	day\_of\_week & TINYINT & Да & Номер дня недели \\ \hline
	start\_time & TIME & Да & Время начала рабочего интервала \\ \hline
	end\_time & TIME & Да & Время окончания рабочего интервала \\ \hline
	weeks\_valid & INT & Нет & Количество недель действия расписания \\ \hline
	is\_active & TINYINT(1) & Нет & Признак активности записи расписания \\ \hline
	created\_at & TIMESTAMP & Да & Дата и время создания записи \\ \hline
\end{longtable}

В состав таблицы service\_unavailable\_dates входят атрибуты, представленные в таблице \ref{unavailabledates:table}.

\begin{longtable}{|>{\raggedright\arraybackslash}p{3.8cm}|>{\raggedright\arraybackslash}p{3.8cm}|>{\centering\arraybackslash}p{2.0cm}|>{\raggedright\arraybackslash}p{5.4cm}|}
	\caption{Атрибуты таблицы service\_unavailable\_dates}\label{unavailabledates:table} \\ \hline
	\centering Поле & \centering Тип & \centering Обязат. & \centering\arraybackslash Описание \\ \hline
	\endfirsthead
	\multicolumn{4}{l}{Продолжение таблицы \ref{unavailabledates:table}} \\ \hline
	\centering Поле & \centering Тип & \centering Обязат. & \centering\arraybackslash Описание \\ \hline
	\endhead
	id & INT & Да & Уникальный идентификатор недоступной даты \\ \hline
	service\_id & INT & Да & Ссылка на услугу \\ \hline
	unavailable\_date & DATE & Да & Дата, исключённая из расписания \\ \hline
	reason & VARCHAR(255) & Нет & Причина недоступности даты \\ \hline
\end{longtable}

В состав таблицы slots входят атрибуты, представленные в таблице \ref{slots:table}.

\begin{longtable}{|>{\raggedright\arraybackslash}p{3.8cm}|>{\raggedright\arraybackslash}p{3.8cm}|>{\centering\arraybackslash}p{2.0cm}|>{\raggedright\arraybackslash}p{5.4cm}|}
	\caption{Атрибуты таблицы slots}\label{slots:table} \\ \hline
	\centering Поле & \centering Тип & \centering Обязат. & \centering\arraybackslash Описание \\ \hline
	\endfirsthead
	\multicolumn{4}{l}{Продолжение таблицы \ref{slots:table}} \\ \hline
	\centering Поле & \centering Тип & \centering Обязат. & \centering\arraybackslash Описание \\ \hline
	\endhead
	id & INT & Да & Уникальный идентификатор слота \\ \hline
	service\_id & INT & Да & Ссылка на услугу \\ \hline
	start\_time & DATETIME & Да & Дата и время начала слота \\ \hline
	status & ENUM(available, booked, unavailable) & Да & Состояние слота \\ \hline
	created\_at & TIMESTAMP & Да & Дата и время создания слота \\ \hline
\end{longtable}

В состав таблицы bookings входят атрибуты, представленные в таблице \ref{bookings:table}.

\begin{longtable}{|>{\raggedright\arraybackslash}p{3.8cm}|>{\raggedright\arraybackslash}p{4.0cm}|>{\centering\arraybackslash}p{2.0cm}|>{\raggedright\arraybackslash}p{5.2cm}|}
	\caption{Атрибуты таблицы bookings}\label{bookings:table} \\ \hline
	\centering Поле & \centering Тип & \centering Обязат. & \centering\arraybackslash Описание \\ \hline
	\endfirsthead
	\multicolumn{4}{l}{Продолжение таблицы \ref{bookings:table}} \\ \hline
	\centering Поле & \centering Тип & \centering Обязат. & \centering\arraybackslash Описание \\ \hline
	\endhead
	id & INT & Да & Уникальный идентификатор бронирования \\ \hline
	service\_id & INT & Да & Ссылка на выбранную услугу \\ \hline
	customer\_id & INT & Да & Ссылка на заказчика \\ \hline
	slot\_id & INT & Нет & Ссылка на выбранный слот \\ \hline
	booking\_date & DATE & Нет & Дата оказания услуги \\ \hline
	price\_at\_booking & DECIMAL(10,2) & Да & Стоимость услуги на момент бронирования \\ \hline
	status & ENUM(pending, confirmed, completed, cancelled, missed) & Да & Текущий статус заявки \\ \hline
	cancelled\_by & ENUM(customer, provider, system) & Нет & Сторона, отменившая заявку \\ \hline
	cancel\_reason & VARCHAR(255) & Нет & Причина отмены заявки \\ \hline
	customer\_message & TEXT & Нет & Сообщение заказчика при оформлении заявки \\ \hline
	booked\_at & TIMESTAMP & Да & Дата и время создания бронирования \\ \hline
	confirmed\_at & TIMESTAMP & Нет & Дата и время подтверждения заявки \\ \hline
	completed\_at & TIMESTAMP & Нет & Дата и время завершения услуги \\ \hline
	updated\_at & TIMESTAMP & Да & Дата и время последнего обновления заявки \\ \hline
\end{longtable}

В состав таблицы chats входят атрибуты, представленные в таблице \ref{chats:table}.

\begin{longtable}{|>{\raggedright\arraybackslash}p{3.8cm}|>{\raggedright\arraybackslash}p{3.8cm}|>{\centering\arraybackslash}p{2.0cm}|>{\raggedright\arraybackslash}p{5.4cm}|}
	\caption{Атрибуты таблицы chats}\label{chats:table} \\ \hline
	\centering Поле & \centering Тип & \centering Обязат. & \centering\arraybackslash Описание \\ \hline
	\endfirsthead
	\multicolumn{4}{l}{Продолжение таблицы \ref{chats:table}} \\ \hline
	\centering Поле & \centering Тип & \centering Обязат. & \centering\arraybackslash Описание \\ \hline
	\endhead
	id & INT & Да & Уникальный идентификатор чата \\ \hline
	booking\_id & INT & Да & Ссылка на бронирование \\ \hline
	created\_at & TIMESTAMP & Да & Дата и время создания чата \\ \hline
\end{longtable}

В состав таблицы messages входят атрибуты, представленные в таблице \ref{messages:table}.

\begin{longtable}{|>{\raggedright\arraybackslash}p{3.8cm}|>{\raggedright\arraybackslash}p{3.8cm}|>{\centering\arraybackslash}p{2.0cm}|>{\raggedright\arraybackslash}p{5.4cm}|}
	\caption{Атрибуты таблицы messages}\label{messages:table} \\ \hline
	\centering Поле & \centering Тип & \centering Обязат. & \centering\arraybackslash Описание \\ \hline
	\endfirsthead
	\multicolumn{4}{l}{Продолжение таблицы \ref{messages:table}} \\ \hline
	\centering Поле & \centering Тип & \centering Обязат. & \centering\arraybackslash Описание \\ \hline
	\endhead
	id & INT & Да & Уникальный идентификатор сообщения \\ \hline
	chat\_id & INT & Да & Ссылка на чат по заявке \\ \hline
	sender\_id & INT & Да & Ссылка на отправителя сообщения \\ \hline
	message & TEXT & Нет & Текст сообщения \\ \hline
	is\_system & TINYINT(1) & Нет & Признак системного сообщения \\ \hline
	is\_read & TINYINT(1) & Нет & Признак прочтения сообщения \\ \hline
	created\_at & TIMESTAMP & Да & Дата и время отправки сообщения \\ \hline
\end{longtable}

В состав таблицы message\_images входят атрибуты, представленные в таблице \ref{messageimages:table}.

\begin{longtable}{|>{\raggedright\arraybackslash}p{3.8cm}|>{\raggedright\arraybackslash}p{3.8cm}|>{\centering\arraybackslash}p{2.0cm}|>{\raggedright\arraybackslash}p{5.4cm}|}
	\caption{Атрибуты таблицы message\_images}\label{messageimages:table} \\ \hline
	\centering Поле & \centering Тип & \centering Обязат. & \centering\arraybackslash Описание \\ \hline
	\endfirsthead
	\multicolumn{4}{l}{Продолжение таблицы \ref{messageimages:table}} \\ \hline
	\centering Поле & \centering Тип & \centering Обязат. & \centering\arraybackslash Описание \\ \hline
	\endhead
	id & INT & Да & Уникальный идентификатор изображения сообщения \\ \hline
	message\_id & INT & Да & Ссылка на сообщение \\ \hline
	image\_path & VARCHAR(255) & Да & Путь к изображению, прикреплённому к сообщению \\ \hline
\end{longtable}

В состав таблицы reviews входят атрибуты, представленные в таблице \ref{reviews:table}.

\begin{longtable}{|>{\raggedright\arraybackslash}p{3.8cm}|>{\raggedright\arraybackslash}p{3.8cm}|>{\centering\arraybackslash}p{2.0cm}|>{\raggedright\arraybackslash}p{5.4cm}|}
	\caption{Атрибуты таблицы reviews}\label{reviews:table} \\ \hline
	\centering Поле & \centering Тип & \centering Обязат. & \centering\arraybackslash Описание \\ \hline
	\endfirsthead
	\multicolumn{4}{l}{Продолжение таблицы \ref{reviews:table}} \\ \hline
	\centering Поле & \centering Тип & \centering Обязат. & \centering\arraybackslash Описание \\ \hline
	\endhead
	id & INT & Да & Уникальный идентификатор отзыва \\ \hline
	booking\_id & INT & Да & Ссылка на бронирование \\ \hline
	author\_id & INT & Да & Ссылка на автора отзыва \\ \hline
	target\_id & INT & Да & Ссылка на пользователя, которому адресован отзыв \\ \hline
	target\_type & ENUM(provider, customer) & Да & Тип адресата отзыва \\ \hline
	rating & TINYINT & Да & Числовая оценка пользователя \\ \hline
	comment & TEXT & Нет & Текст комментария к отзыву \\ \hline
	created\_at & TIMESTAMP & Да & Дата и время создания отзыва \\ \hline
\end{longtable}

В состав таблицы favorites входят атрибуты, представленные в таблице \ref{favorites:table}.

\begin{longtable}{|>{\raggedright\arraybackslash}p{3.8cm}|>{\raggedright\arraybackslash}p{3.8cm}|>{\centering\arraybackslash}p{2.0cm}|>{\raggedright\arraybackslash}p{5.4cm}|}
	\caption{Атрибуты таблицы favorites}\label{favorites:table} \\ \hline
	\centering Поле & \centering Тип & \centering Обязат. & \centering\arraybackslash Описание \\ \hline
	\endfirsthead
	\multicolumn{4}{l}{Продолжение таблицы \ref{favorites:table}} \\ \hline
	\centering Поле & \centering Тип & \centering Обязат. & \centering\arraybackslash Описание \\ \hline
	\endhead
	id & INT & Да & Уникальный идентификатор записи избранного \\ \hline
	user\_id & INT & Да & Ссылка на пользователя \\ \hline
	service\_id & INT & Да & Ссылка на услугу, добавленную в избранное \\ \hline
	created\_at & TIMESTAMP & Да & Дата и время добавления услуги в избранное \\ \hline
\end{longtable}

Связи между сущностями реализуются с использованием внешних ключей. Пользователь-исполнитель может быть связан с несколькими услугами, каждая услуга может иметь множество расписаний, недоступных дат, слотов и бронирований, а каждое бронирование связано с одним заказчиком и одной услугой. Для каждого бронирования создаётся отдельный чат, содержащий сообщения, которые могут сопровождаться изображениями. После завершения услуги каждая сторона может оставить отзыв о другой стороне, а для предотвращения повторного добавления используется ограничение уникальности на уровне таблицы reviews.